\svnid{$Id: introduction.tex 74563 2022-02-13 20:35:17Z jagers $}

\chapter{Introduction}
The MATLAB interface for Delft3D has been developed to expand the flexibility in the processing of Delft3D data in general and of NEFIS data files in particular.
The interface consists of three parts:

\begin{itemize}
\item First, generic low level routines to access data stored in NEFIS files by the \dfoursuite.
The routines are described in the second chapter of this manual.
A thorough knowledge of the Delft3D system is required.

\item Second, high level routines to access the data.
This includes routines for reading and writing \dflow restart files, grid and depth files.
Routines to load and plot grid, thin dams, and velocity data directly from the \dfoursuite result files.
Finally a set of routines to open and read data from any file supported by QUICKPLOT; this includes the result files of both the \dfoursuite and the \dfmsuite.
For a full overview of all supported file formats see the appendix of \citet{QP_UM}.
These routines are described in Chapter \ref{Chp:HighLevelRoutines}.

\item Third, a graphical user interface for basic plotting.
The QUICKPLOT\ interface is described in Chapter \ref{Chp:QuickPlot}.
\end{itemize}

All functions are provided in source code under LGPL v2.1 conditions.
The routines described in chapters three and four can be used as examples for development of other routines.

\section{Version information}
This manual describes the functionality of \DMATLAB\ interface version 2.60 and later versions.

\section{Compatibility}
The routines have been tested for compatibility with MATLAB versions ranging from 6.5 (R13) until 2022a.
Although quotes in commands are sometimes rendered in this manual as \command{'}, every quote is intended to be a straight quote \textcode{'}.

